\documentclass[12pt, a4paper]{article}

% ── PACKAGES ──────────────────────────────────────────────────
\usepackage[utf8]{inputenc}
\usepackage[T1]{fontenc}
\usepackage{geometry}
\geometry{margin=2.5cm}
\usepackage{hyperref}
\hypersetup{
    colorlinks=true,
    linkcolor=blue,
    urlcolor=blue,
    citecolor=blue,
    pdftitle={Fulvestrant Superiority Over Aromatase
              Inhibitors in FOXA1-Hyperactivated Invasive
              Lobular Carcinoma: A Geometry-Derived
              Mechanistic Prediction with Proposed
              Validation in NCT02206984},
    pdfauthor={Eric Robert Lawson},
    pdfkeywords={ILC, invasive lobular carcinoma,
                 fulvestrant, aromatase inhibitor, FOXA1,
                 CDH1, E-cadherin, ER circuit amplification,
                 endocrine therapy, SERD, breast cancer,
                 NCT02206984, FOXA1-stratified, Type 3,
                 structural lock, attractor geometry,
                 Waddington landscape, OrganismCore,
                 GSE176078, METABRIC, treatment prediction,
                 precision oncology, ILC endocrine resistance}
}
\usepackage{amsmath}
\usepackage{booktabs}
\usepackage{longtable}
\usepackage{array}
\usepackage{parskip}
\usepackage{microtype}
\usepackage{authblk}
\usepackage{titlesec}
\usepackage{fancyhdr}
\usepackage{xcolor}
\usepackage{float}
\usepackage{enumitem}
\usepackage{setspace}
\usepackage{multirow}

% ── COLOURS ───────────────────────────────────────────────────
\definecolor{repobox}{RGB}{240,247,255}
\definecolor{findingbox}{RGB}{245,255,245}
\definecolor{convergentbox}{RGB}{255,250,235}
\definecolor{mechanismbox}{RGB}{250,245,255}
\definecolor{novelbox}{RGB}{255,248,220}
\definecolor{actionbox}{RGB}{230,255,230}
\definecolor{urgentbox}{RGB}{255,235,235}

% ── HEADER / FOOTER ───────────────────────────────────────────
\pagestyle{fancy}
\fancyhf{}
\lhead{\footnotesize CS-LIT-18: Fulvestrant $>$ AI
       in FOXA1-Stratified ILC}
\rhead{\footnotesize OrganismCore \textbar{} 2026-03-06}
\rfoot{\small Page \thepage}
\renewcommand{\headrulewidth}{0.4pt}
\setlength{\headheight}{24pt}
\addtolength{\topmargin}{-10pt}

% ── SECTION FORMATTING ────────────────────────────────────────
\titleformat{\section}
  {\large\bfseries}{\thesection}{1em}{}[\titlerule]
\titleformat{\subsection}
  {\normalsize\bfseries}{\thesubsection}{1em}{}
\titleformat{\subsubsection}
  {\small\bfseries}{\thesubsubsection}{1em}{}

% ══════════════════════════════════════════════════════════════
% TITLE
% ══════════════════════════════════════════════════════════════
\title{
    \vspace{-1em}
    \textbf{Fulvestrant Superiority Over Aromatase
    Inhibitors in FOXA1-Hyperactivated Invasive
    Lobular Carcinoma:}\\[0.5em]
    \large A Geometry-Derived Mechanistic Prediction
    with Proposed Validation in NCT02206984 and
    Archived ILC Tissue Cohorts\\[0.4em]
    \normalsize \textit{OrganismCore Framework ---
    Document CS-LIT-18}\\[0.3em]
    \small Derived from Waddington Attractor Geometry
    of 19,542 Single Cancer Cells (GSE176078)\\[0.3em]
    \small Companion to CS-LIT-7: ILC as the Geometric
    Inverse of TNBC
    (\href{https://doi.org/10.5281/zenodo.18893825}
    {DOI: 10.5281/zenodo.18893825})
}

\author[1]{Eric Robert Lawson}
\affil[1]{
    OrganismCore\\
    \href{mailto:OrganismCore@proton.me}
         {OrganismCore@proton.me}\\
    ORCID:
    \href{https://orcid.org/0009-0002-0414-6544}
         {0009-0002-0414-6544}
}

\date{March 6, 2026}

% ══════════════════════════════════════════════════════════════
\begin{document}
% ══════════════════════════════════════════════════════════════

\maketitle
\thispagestyle{fancy}

% ── REPOSITORY POINTER ────────────────────────────────────────
\begin{center}
\colorbox{repobox}{
\begin{minipage}{0.88\textwidth}
\vspace{0.5em}
\centering
\textbf{Full computational record --- scripts, data
caches, reasoning artifacts, and raw logs:}\\[0.4em]
\href{https://github.com/Eric-Robert-Lawson/attractor-oncology}
{\texttt{https://github.com/Eric-Robert-Lawson/attractor-oncology}}
\\[0.3em]
\small Repository ID: 1172048282 \textbar{}
MIT License \textbar{}
Python \textbar{}
Public\\[0.2em]
Part of the OrganismCore BRCA Cross-Subtype Analysis
(BRCA-S8h, 2026-03-05).\\
All predictions locked from attractor geometry
\textbf{before} any literature review was performed.
\vspace{0.5em}
\end{minipage}}
\end{center}

\vspace{0.8em}

% ── URGENCY BOX ───────────────────────────────────────────────
\begin{center}
\colorbox{urgentbox}{
\begin{minipage}{0.88\textwidth}
\vspace{0.5em}
\centering
\textbf{ACTIVE TRIAL NOTE}\\[0.4em]
\small
NCT02206984 --- Endocrine Response in Women With
Invasive Lobular Breast Cancer (pre-surgical window
study; fulvestrant vs AI vs tamoxifen; Ki67 endpoint;
University of Pittsburgh / UPMC).\\[0.3em]
PI: Dr.\ Priscilla McAuliffe, MD PhD, UPMC.\\
Trial contact: Brenda Steele (steebx@upmc.edu).\\[0.3em]
This trial does \textbf{not} currently use FOXA1
expression as a patient stratification or correlative
variable. The FOXA1 H-score on pre-treatment biopsy
tissue --- already obtained as part of the pre-surgical
design --- could be added as a correlative endpoint
without protocol amendment or additional patient burden.\\[0.3em]
Companion ILC researchers: Dr.\ Adrian Lee and
Dr.\ Steffi Oesterreich (Lee/Oesterreich Laboratory,
Institute for Precision Medicine, UPMC Hillman Cancer
Center) --- the leading ILC molecular research group
in North America.
\vspace{0.5em}
\end{minipage}}
\end{center}

\vspace{0.8em}

% ══════════════════════════════════════════════════════════════
\begin{abstract}
% ══════════════════════════════════════════════════════════════
\noindent
Invasive Lobular Carcinoma (ILC) is the second most
common breast cancer histology (${\sim}10$--$15\%$ of
all cases) and is universally ER-positive, yet it is
treated with the same endocrine agents as Invasive
Ductal Carcinoma (IDC). The OrganismCore Waddington
attractor geometry framework identifies ILC as a
\textbf{Type~3 structural lock}: FOXA1 is hyperactivated
above normal Luminal A levels, CDH1 (E-cadherin) is
absent (${\sim}65\%$ mutation, ${\sim}35\%$ methylation),
and EZH2 is low. FOXA1 is the pioneer transcription
factor that controls ER circuit gain --- it physically
opens chromatin for ER binding and amplifies ER
transcriptional output. Hyperactivation of FOXA1 in
ILC means the ER circuit operates at above-normal
gain.

This single geometrical observation generates a
specific, testable, and falsifiable therapeutic
prediction: in FOXA1-hyperactivated ILC, aromatase
inhibitors (AIs) are mechanistically insufficient
as the primary endocrine strategy because they act
by reducing ligand (estrogen), and the FOXA1-amplified
ER circuit partially compensates for reduced ligand
by maintaining chromatin accessibility and
transcriptional readiness at ER target loci. Full
ER receptor degradation (fulvestrant, a selective
ER degrader/SERD) eliminates the receptor protein
upstream of FOXA1 amplification, producing complete
suppression of ER circuit output regardless of FOXA1
gain.

The prediction: \textbf{in FOXA1-high ILC,
fulvestrant produces superior PFS and/or Ki67
suppression compared to AIs}, with the advantage
increasing as a function of FOXA1 expression level.
This is not a prediction based on ESR1 mutation
(the conventional rationale for fulvestrant
preference) but on FOXA1 circuit amplification
as a class property of ILC.

The biological rationale is confirmed by published
ILC biology: Frontiers Oncology 2025 ILC review
explicitly notes FOXA1 hyperactivation and altered
endocrine sensitivity as an open research question
directly adjacent to this prediction. The individual
features (FOXA1-high ILC, CDH1-absent ILC, FOXA1 as
ER circuit amplifier) are each published. The
mechanistic derivation that FOXA1 amplification
specifically undermines AI efficacy and specifically
favours fulvestrant --- and the patient selection
criterion based on FOXA1 IHC H-score --- is not
published in any trial design or mechanistic paper.
Verdict: CONVERGENT-NOVEL.

NCT02206984 (pre-surgical endocrine comparison in
ILC, UPMC, Ki67 endpoint, PI: Dr.\ Priscilla
McAuliffe) is the immediate clinical context for
this prediction. FOXA1 IHC on pre-treatment biopsy
tissue --- already obtained as part of the
pre-surgical design --- could be added as a
correlative endpoint without protocol amendment.
\end{abstract}

\vspace{0.5em}
\noindent\textbf{Keywords:}
ILC; invasive lobular carcinoma; fulvestrant; SERD;
aromatase inhibitor; FOXA1; CDH1; E-cadherin; ER
circuit amplification; endocrine therapy; breast cancer;
NCT02206984; FOXA1-stratified; Type~3; structural lock;
attractor geometry; Waddington landscape; OrganismCore;
GSE176078; treatment prediction; precision oncology

\newpage
\tableofcontents
\newpage

% ══════════════════════════════════════════════════════════════
\section{Document Metadata}
% ══════════════════════════════════════════════════════════════

\begin{center}
\begin{tabular}{ll}
\toprule
\textbf{Field} & \textbf{Value} \\
\midrule
Document ID & CS-LIT-18 \\
Parent document & BRCA-S8h (2026-03-05) \\
Verdict & CONVERGENT-NOVEL \\
Date & 2026-03-06 \\
Author & Eric Robert Lawson \\
ORCID & 0009-0002-0414-6544 \\
Organisation & OrganismCore \\
Contact & OrganismCore@proton.me \\
Repository &
\href{https://github.com/Eric-Robert-Lawson/attractor-oncology}
{attractor-oncology} (ID: 1172048282) \\
Derivation data & GSE176078 (19,542 single cancer cells) \\
Requires & CS-LIT-7 (foundational architecture) \\
CS-LIT-7 DOI &
\href{https://doi.org/10.5281/zenodo.18893825}
{10.5281/zenodo.18893825} \\
CS-LIT-1 DOI &
\href{https://doi.org/10.5281/zenodo.18883922}
{10.5281/zenodo.18883922} \\
CS-LIT-2 DOI &
\href{https://doi.org/10.5281/zenodo.18884158}
{10.5281/zenodo.18884158} \\
Target trial & NCT02206984 (UPMC, PI: McAuliffe) \\
Trial contact & steebx@upmc.edu (Brenda Steele) \\
License & CC BY 4.0 \\
Biological contradictions & 0 \\
\bottomrule
\end{tabular}
\end{center}

% ══════════════════════════════════════════════════════════════
\section{Background: ILC Endocrine Treatment
and the Unsolved Problem}
% ══════════════════════════════════════════════════════════════

\subsection{What is currently known}

ILC is treated with the same endocrine agents as
IDC: aromatase inhibitors (anastrozole, letrozole,
exemestane), tamoxifen, and fulvestrant. There is
no current ILC-specific endocrine treatment
algorithm. Guidelines treat ER-positive breast
cancer as a single class for endocrine agent
selection.

Published clinical data on ILC vs IDC endocrine
outcomes:

\begin{itemize}
    \item Several retrospective analyses suggest
    AIs may be less effective in ILC than in IDC
    relative to tamoxifen (Cristofanilli et al.,
    Moulie et al., others).
    \item A meta-analysis of ATAC and BIG 1-98
    subgroup data suggested differential AI vs
    tamoxifen benefit by histology, but results
    have been inconsistent across studies.
    \item The Frontiers Oncology 2025 ILC review
    explicitly identifies FOXA1 hyperactivation
    and altered endocrine sensitivity as an open
    research question.
    \item NCT02206984 (UPMC) compares fulvestrant,
    AI (letrozole), and tamoxifen in the
    pre-surgical window in ILC using Ki67 as the
    primary endpoint --- but does \textit{not}
    stratify by FOXA1 expression.
\end{itemize}

\subsection{The gap}

No published trial or analysis has:
\begin{enumerate}
    \item Stratified endocrine agent choice in ILC
    by FOXA1 expression level
    \item Provided a mechanistic explanation for
    why FOXA1 hyperactivation specifically undermines
    AI efficacy (as opposed to ESR1 mutation, which
    is the conventional mechanism for fulvestrant
    preference)
    \item Generated a falsifiable prediction that
    FOXA1 IHC H-score is a continuous predictor of
    fulvestrant benefit magnitude over AI in ILC
\end{enumerate}

The OrganismCore framework provides all three.

% ══════════════════════════════════════════════════════════════
\section{The Geometric Basis: ILC as a Type~3
Structural Lock}
% ══════════════════════════════════════════════════════════════

This section summarises the architectural basis
established in CS-LIT-7 (ILC as the geometric
inverse of TNBC,
\href{https://doi.org/10.5281/zenodo.18893825}
{DOI: 10.5281/zenodo.18893825}).

\subsection{ILC on the luminal identity axis}

The FOXA1/EZH2 ratio orders all breast cancer
subtypes on a continuous luminal identity axis
(CS-LIT-1,
\href{https://doi.org/10.5281/zenodo.18883922}
{DOI: 10.5281/zenodo.18883922}). ILC sits
\textit{above} LumA on this axis --- FOXA1 is
hyperactivated beyond the normal luminal maximum.
The axis runs:

\begin{center}
\small
CL ($0.10$) $\leftarrow$
TNBC ($0.52$) $\leftarrow$
HER2 ($3.34$) $\leftarrow$
LumB ($8.10$) $\leftarrow$
LumA ($9.38$) $\leftarrow$
\textbf{ILC (FOXA1 $\uparrow\uparrow$ above LumA)}
\end{center}

\subsection{What FOXA1 hyperactivation means
for ER circuit gain}

FOXA1 is the pioneer transcription factor that
defines luminal identity. Its specific function:
it physically remodels compacted chromatin at
ER-target loci, creating accessible nucleosome-free
regions into which the ER can bind. FOXA1 does not
merely permit ER binding --- it actively amplifies
ER transcriptional output by:

\begin{itemize}
    \item Maintaining chromatin accessibility at
    ER-bound loci above baseline (the ER finds
    pre-opened chromatin rather than having to
    compete with nucleosomes)
    \item Stabilising ER occupancy at binding
    sites through FOXA1/ER co-occupancy
    \item Recruiting co-activators to the same
    loci, amplifying the transcriptional response
    per ER binding event
\end{itemize}

In a normal LumA cell, FOXA1 operates at a
calibrated level: the ER circuit has defined gain.
In ILC, FOXA1 is overexpressed beyond this
calibration. The ER circuit operates at elevated
gain --- the same amount of ER activity produces
more transcriptional output than it would in a
LumA cell.

\subsection{CDH1 loss: the structural component}

CDH1 (E-cadherin) loss in ILC
(${\sim}65\%$ mutation, ${\sim}35\%$ methylation)
removes the structural adhesion constraint.
This is the Type~3 classification: the lock is
structural, not epigenetic. The identity programme
(FOXA1 hyperactive, ER circuit amplified) is intact
and overdriven. The cell is disinhibited structurally
while being over-committed to luminal identity
molecularly.

% ══════════════════════════════════════════════════════════════
\section{The Therapeutic Prediction: Mechanism}
% ══════════════════════════════════════════════════════════════

\subsection{How aromatase inhibitors work in
normal luminal cells}

Aromatase inhibitors block the conversion of
androgens to estrogens in peripheral tissues,
reducing circulating estrogen to near-zero in
postmenopausal women. In a normal luminal cell:

\begin{enumerate}
    \item Estrogen binds ER
    \item ER-estrogen complex recruits FOXA1
    at EREs (oestrogen response elements)
    \item ER-FOXA1 complex drives target
    gene transcription (TFF1, GATA3, PGR, etc.)
    \item AI reduces step~1 input: less estrogen
    $\rightarrow$ less ER activation $\rightarrow$
    less transcriptional output
\end{enumerate}

In a normal LumA cell, this chain is approximately
linear: halving ligand roughly halves ER circuit
output. AI efficacy depends on this linearity.

\subsection{Why FOXA1 hyperactivation breaks
AI linearity in ILC}

In an ILC cell with FOXA1 hyperactivated:

\begin{enumerate}
    \item FOXA1 maintains persistent chromatin
    accessibility at ER target loci independent
    of estrogen level
    \item Even with reduced estrogen (AI effect),
    the ER-target chromatin is pre-opened and
    co-activator complexes are pre-positioned
    \item The reduced estrogen signal activates
    a FOXA1-primed circuit rather than an
    unprimed one --- each remaining ER binding
    event produces more output than in LumA
    \item The AI suppresses input but the
    FOXA1-amplified circuit partially compensates
    through increased output per activation event
\end{enumerate}

\begin{center}
\colorbox{mechanismbox}{
\begin{minipage}{0.88\textwidth}
\vspace{0.5em}
\textbf{The gain amplification argument:}\\[0.4em]
If an AI reduces estrogen input by $80\%$, and
the LumA circuit gain is $\times1$, ER output
is reduced to $20\%$.\\[0.3em]
If an AI reduces estrogen input by $80\%$, and
the ILC circuit gain is $\times3$ (FOXA1
hyperactivation), ER output is reduced to
${\sim}60\%$ (the elevated gain partially
compensates for reduced input).\\[0.3em]
The AI achieves the same ligand suppression.
The ILC cell has higher residual ER circuit
output because FOXA1 amplification operates
independently of ligand level.
\vspace{0.5em}
\end{minipage}}
\end{center}

\subsection{Why fulvestrant overcomes this}

Fulvestrant is a selective ER degrader (SERD).
Its mechanism:

\begin{enumerate}
    \item Fulvestrant binds the ER protein with
    high affinity
    \item The fulvestrant-ER complex is targeted
    for proteasomal degradation
    \item ER protein is eliminated from the cell
    \item With no ER protein, FOXA1 amplification
    is irrelevant: there is no receptor to amplify
    \item ER circuit output drops to zero
    independent of FOXA1 level
\end{enumerate}

Fulvestrant acts upstream of FOXA1 in the
signalling chain. It does not compete with FOXA1
amplification --- it eliminates the object that
FOXA1 was amplifying. This is why fulvestrant
efficacy is predicted to be FOXA1-independent
in ILC (it works regardless of FOXA1 level) while
AI efficacy is predicted to be FOXA1-dependent
(attenuated by higher FOXA1).

\subsection{The testable prediction stated formally}

\begin{center}
\colorbox{findingbox}{
\begin{minipage}{0.88\textwidth}
\vspace{0.5em}
\textbf{THE FORMAL PREDICTION:}\\[0.4em]
In ILC patients, \textbf{fulvestrant produces
greater Ki67 suppression, PFS benefit, and/or OS
benefit compared to AIs}, and the \textbf{magnitude
of this advantage increases as a continuous
function of FOXA1 IHC H-score at baseline}.\\[0.4em]
\textbf{Corollary 1:} In FOXA1-low ILC (if any
such cases exist with classical ILC morphology),
the fulvestrant vs AI advantage will be
attenuated or absent.\\[0.4em]
\textbf{Corollary 2:} The fulvestrant advantage
in ILC is not primarily ESR1-mutation-dependent
--- it is present in ESR1 wild-type FOXA1-high
ILC. This distinguishes the geometric prediction
from the conventional rationale.\\[0.4em]
\textbf{Falsification:} FOXA1 H-score does not
correlate with differential Ki67 suppression
(fulvestrant vs AI) in NCT02206984 tissue,
or shows negative/null correlation.
\vspace{0.5em}
\end{minipage}}
\end{center}

% ══════════════════════════════════════════════════════════════
\section{What Is and Is Not in the Published
Literature}
% ══════════════════════════════════════════════════════════════

\subsection{What IS established}

\begin{enumerate}
    \item \textbf{FOXA1 hyperactivation in ILC:}
    Published. Frontiers Oncology 2025 ILC review
    notes FOXA1 hyperactivation in ILC and
    identifies altered endocrine sensitivity as
    an open question. FOXA1 gain-of-function
    mutations in ILC published (Cancer Cell 2020).

    \item \textbf{FOXA1 as ER circuit amplifier:}
    Published. FOXA1 pioneer factor function,
    chromatin opening at EREs, and co-activation
    of ER-driven transcription are well-established
    molecular biology.

    \item \textbf{Differential AI vs tamoxifen
    in ILC vs IDC:}
    Multiple retrospective analyses suggest
    differential effects. Results are inconsistent
    but the field acknowledges ILC may respond
    differently to endocrine agents.

    \item \textbf{Fulvestrant mechanism (SERD):}
    Established. Fulvestrant degrades ER protein.
    FDA-approved. Active in post-AI endocrine
    resistance.

    \item \textbf{NCT02206984 (active):}
    Pre-surgical comparison of fulvestrant, AI,
    and tamoxifen in ILC. Ki67 endpoint. UPMC.
    PI: Dr.\ Priscilla McAuliffe. Does not include
    FOXA1 stratification.
\end{enumerate}

\subsection{What is NOT in the literature}

\begin{center}
\colorbox{novelbox}{
\begin{minipage}{0.88\textwidth}
\vspace{0.5em}
\textbf{The following have no published equivalent
as of 2026-03-06:}
\begin{itemize}[noitemsep]
    \item The mechanistic derivation that FOXA1
    hyperactivation in ILC specifically attenuates
    AI efficacy by maintaining ER circuit gain
    independent of ligand level.
    \item The prediction that fulvestrant
    superiority over AI in ILC is a FOXA1
    circuit amplification effect (not ESR1
    mutation), present in ESR1 wild-type
    FOXA1-high ILC.
    \item The use of FOXA1 IHC H-score as a
    continuous predictor of the magnitude of
    fulvestrant benefit over AI in ILC.
    \item Any published trial stratifying
    fulvestrant vs AI in ILC by FOXA1 expression.
    \item The falsifiable prediction that the
    advantage increases as a continuous function
    of FOXA1 level.
\end{itemize}
\vspace{0.5em}
\end{minipage}}
\end{center}

% ══════════════════════════════════════════════════════════════
\section{Clinical Validation Pathway}
% ══════════════════════════════════════════════════════════════

\subsection{NCT02206984: Immediate opportunity}

\begin{center}
\colorbox{actionbox}{
\begin{minipage}{0.88\textwidth}
\vspace{0.5em}
\textbf{NCT02206984 TRIAL DETAILS:}\\[0.4em]
\textbf{Title:} Endocrine Response in Women With
Invasive Lobular Breast Cancer (pre-surgical
window study).\\[0.3em]
\textbf{Sponsor:} UPMC / University of Pittsburgh.\\[0.3em]
\textbf{PI:} Dr.\ Priscilla McAuliffe, MD PhD.\\[0.3em]
\textbf{Trial contact:} Brenda Steele
(steebx@upmc.edu) / Rometa Pollard
(pollardrr@upmc.edu).\\[0.3em]
\textbf{Arms:} Fulvestrant vs AI (letrozole) vs
tamoxifen in pre-surgical window in ILC.\\[0.3em]
\textbf{Primary endpoint:} Ki67 suppression.\\[0.3em]
\textbf{Tissue:} Pre-treatment core biopsy tissue
already obtained as part of the design.\\[0.4em]
\textbf{The proposed addition:}\\
FOXA1 IHC H-score on pre-treatment core biopsy
as a correlative variable. No additional patient
visit, no additional biopsy, no protocol amendment
to primary endpoint. One additional IHC stain on
already-banked tissue.\\[0.4em]
\textbf{The falsifiable analysis:}\\
Does FOXA1 H-score at baseline correlate with
the magnitude of Ki67 suppression advantage for
fulvestrant over AI? Is the correlation continuous
(higher FOXA1 $\rightarrow$ greater fulvestrant
advantage)?
\vspace{0.5em}
\end{minipage}}
\end{center}

\subsection{ILC tissue archive analysis}

Beyond NCT02206984, the prediction is testable
in any archived ILC tissue cohort with:
\begin{itemize}
    \item FFPE tissue from primary ILC
    \item Endocrine treatment data (AI vs fulvestrant)
    \item Clinical outcome data (PFS, OS, or Ki67)
    \item Available for FOXA1 IHC staining
\end{itemize}

Candidate cohorts include the UPMC ILC Biorepository
(Lee/Oesterreich Laboratory), the ILC-specific
arms of METABRIC (if endocrine agent data are
available), and archived tissue from FALCON
(fulvestrant vs anastrozole in advanced HR$+$
breast cancer; if ILC subgroup data are accessible).

\subsection{Proposed trial design for prospective
validation}

A dedicated FOXA1-stratified ILC endocrine trial
would require:

\begin{center}
\begin{tabular}{p{3.5cm}p{8cm}}
\toprule
\textbf{Parameter} & \textbf{Specification} \\
\midrule
Population &
ER$+$ ILC (classical histology confirmed), any
stage, initiating or changing endocrine therapy \\[0.4em]
Stratification &
FOXA1 IHC H-score on pre-treatment biopsy:
high ($\geq150$) vs intermediate ($50$--$149$)
vs low ($<50$) \\[0.4em]
Arms &
Fulvestrant $500$~mg q28d vs letrozole $2.5$~mg
daily (or anastrozole) \\[0.4em]
Primary endpoint &
Ki67 suppression at 2--4 weeks (pre-surgical)
or PFS at 12 months (metastatic) \\[0.4em]
Key secondary &
ESR1 mutation status (to disentangle FOXA1
mechanism from ESR1 mechanism) \\[0.4em]
Falsification criterion &
No correlation between FOXA1 H-score and
differential Ki67 suppression (fulvestrant vs AI) \\
\bottomrule
\end{tabular}
\end{center}

% ══════════════════════════════════════════════════════════════
\section{Relation to CS-LIT-7 and the
Published Series}
% ══════════════════════════════════════════════════════════════

CS-LIT-18 is the clinical drug prediction document
that follows directly from CS-LIT-7 (the geometric
architecture). Their relationship:

\begin{center}
\begin{tabular}{p{2.8cm}p{8.5cm}}
\toprule
\textbf{Document} & \textbf{Role} \\
\midrule
CS-LIT-7\\
\href{https://doi.org/10.5281/zenodo.18893825}
{DOI: 10.5281/zenodo.18893825} &
The geometric architecture: ILC as the structural
inverse of TNBC. FOXA1 hyperactivated in ILC.
CDH1 absent. Type~3 structural lock. Required
reading before CS-LIT-18. \\[0.5em]
CS-LIT-18\\(this document) &
The clinical prediction: fulvestrant superiority
over AIs in FOXA1-high ILC, with NCT02206984 as
the immediate testing context. \\[0.5em]
CS-LIT-1\\
\href{https://doi.org/10.5281/zenodo.18883922}
{DOI: 10.5281/zenodo.18883922} &
The FOXA1/EZH2 ratio ordering axis. ILC sits
above LumA (FOXA1 hyperactivated). The ratio
instrument that places ILC correctly in the
treatment classification. \\[0.5em]
CS-LIT-2\\
\href{https://doi.org/10.5281/zenodo.18884158}
{DOI: 10.5281/zenodo.18884158} &
The six lock type classification. ILC is Lock
Type~3 (structural). CS-LIT-18 is the therapeutic
evidence document for this lock type. \\[0.5em]
CS-LIT-22\\
\href{https://doi.org/10.5281/zenodo.18892788}
{DOI: 10.5281/zenodo.18892788} &
FOXA1/EZH2 dual IHC tool. The patient selection
instrument that identifies FOXA1-high ILC in the
pathology laboratory. \\
\bottomrule
\end{tabular}
\end{center}

% ══════════════════════════════════════════════════════════════
\section{Locked Statement}
% ══════════════════════════════════════════════════════════════

\begin{center}
\colorbox{findingbox}{
\begin{minipage}{0.88\textwidth}
\vspace{0.6em}
\textbf{Stated once, precisely, without
inflation:}\\[0.5em]
In invasive lobular carcinoma, FOXA1 is
hyperactivated above normal Luminal~A levels,
amplifying the ER circuit gain. Aromatase
inhibitors reduce estrogen ligand; the
FOXA1-amplified circuit partially compensates
for reduced ligand by maintaining chromatin
accessibility and co-activator positioning at
ER target loci, attenuating AI efficacy.
Fulvestrant degrades the ER receptor upstream
of FOXA1 amplification, eliminating circuit
output independent of FOXA1 level. The prediction
is: in FOXA1-high ILC, fulvestrant produces
superior Ki67 suppression and PFS compared to
AIs, with the advantage increasing continuously
with FOXA1 IHC H-score at baseline. This effect
operates in ESR1 wild-type FOXA1-high ILC and
is mechanistically distinct from the conventional
ESR1-mutation rationale for fulvestrant. NCT02206984
(UPMC; pre-surgical fulvestrant vs AI vs tamoxifen
in ILC; Ki67 endpoint; PI: Dr.\ Priscilla McAuliffe)
is the immediate context for this prediction.
FOXA1 IHC on pre-treatment biopsy tissue can be
added as a correlative variable without protocol
amendment. No published trial has stratified
fulvestrant vs AI in ILC by FOXA1 expression.
The mechanistic prediction has no published
equivalent as of 2026-03-06.
Verdict: CONVERGENT-NOVEL.
\vspace{0.5em}
\end{minipage}}
\end{center}

\end{document}
